% This must be in the first 5 lines to tell arXiv to use pdfLaTeX, which is strongly recommended.
\pdfoutput=1
% In particular, the hyperref package requires pdfLaTeX in order to break URLs across lines.

\documentclass[11pt]{article}

% Remove the "review" option to generate the final version.
% \usepackage[review]{acl}
\usepackage{acl}

% Standard package includes
\usepackage{times}
\usepackage{latexsym}
\usepackage[T1]{fontenc}
\usepackage[utf8]{inputenc}
\usepackage{microtype}

% Useful packages for this paper
\usepackage{booktabs}     % nicer tables
\usepackage{multirow}
\usepackage{amsmath}
\usepackage{amssymb}
\usepackage{graphicx}
\usepackage{xcolor}
\usepackage{xspace}

% Convenience macros — adjust/remove as needed
\newcommand{\bertbase}{BERT-base-uncased\xspace}
\newcommand{\nselected}{n_{\text{sel}}\xspace}
\newcommand{\flip}{\text{flip}\xspace}
\newcommand{\todo}[1]{\textcolor{red}{[TODO: #1]}}

\title{A Mechanistic Study of AI-Detection Features in Frozen BERT:\\
Sparse Probing and Activation Patching on RAID}
% Working title — alternatives in the planning notes; pick at submission time.

\author{Anonymous ARR submission}
% \author{
%   Pawe{\l} \\
%   Affiliation \\
%   \texttt{email@domain} \\
%   \And
%   Mi{\l}osz Grunwald \\
%   Affiliation \\
%   \texttt{email@domain} \\
% }

% =====================================================================
% HARD SUBMISSION CRITERIA — all must be resolved before May 25
% =====================================================================
% [ ] ABSTRACT — write the 6-sentence abstract (S1 context, S2 gap, S3 this paper,
%     S4 key finding, S5 validation, S6 impact). Currently a \todo{} placeholder.
%
% [ ] LIMITATIONS SECTION — mandatory for ARR; does NOT count toward the 8-page limit.
%     Must cover: (1) single backbone (frozen BERT only), (2) generalizability across
%     domains/languages, (3) causal scope (patching is not a full circuit proof),
%     (4) computational scope (5 generators, not all 11 in RAID).
%
% [ ] ALL \todo{} PLACEHOLDERS RESOLVED — search for \todo{ before submission.
%     Currently open: abstract, setup subsections, results subsections, discussion.
%
% [ ] SETUP SECTION FILLED — §3 needs: RAID data description with statistics,
%     model + feature extraction details (BERT-base-uncased, CLS, 12×768),
%     probing protocol (L1→L2, C=0.005, N=7500, 5-fold × 3 seeds),
%     patching protocol (same-domain donor, mean ablation), LOGO protocol.
%
% [ ] RESULTS SECTION FILLED — §4/§5 needs all four experiment types reported:
%     sparse probe (accuracy + Jaccard), patching (flip rates + ratio vs random),
%     ablation (accuracy drop), cross-generator (pairwise Jaccard, layer share),
%     LOGO (sel-L2 accuracy across held-out families).
%
% [ ] MATHEMATICAL FORMALIZATION — at minimum: notation for BERT CLS activations
%     h^{(l)} \in \mathbb{R}^{768}, concatenated representation z \in \mathbb{R}^{9216},
%     L1 probe definition, mean-ablation definition, patching operation definition.
%     Expected in every accepted MI paper; signals rigor to reviewers.
%
% [ ] FIGURE 1 (PIPELINE DIAGRAM) — almost universally present in accepted MI papers.
%     Should show: BERT encoder → CLS extraction → L1 probe → stable neuron set →
%     [mean ablation branch] and [donor patching branch] → flip rate metric.
%
% [ ] HEDGE RESULTS CLAIMS — use "suggest", "provide evidence that", "indicate"
%     rather than "demonstrate" or "prove" in results prose. Check intro ¶4 too.
%
% [ ] REFERENCE COUNT — currently 14 refs; accepted MI papers average 40–70.
%     Add foundational MI refs as sections are written (Elhage 2021 circuits framework,
%     Geva 2021 FFN-as-key-value if we discuss FFN neurons, others as needed).
%
% [ ] ETHICS STATEMENT — optional but expected; one short paragraph is sufficient.
%     Note: detection tools can be misused; frozen BERT study is observational, not a
%     deployed system.
% =====================================================================

\begin{document}
\maketitle

% =====================================================================
\begin{abstract}
% ~150–200 words, single paragraph.
% Recommended structure (one sentence each):
% (1) Problem: AI-text detection is empirically solvable but mechanistically opaque on encoder classifiers.
% (2) Setup: frozen \bertbase on RAID, 5 generators (1 base + 4 chat-RLHF), Gurnee-style L1→L2 sparse probing over all 12×768=9{,}216 CLS features.
% (3) Result A (sparse stable set): 0.5–0.9\% of features (45–62 stable neurons) suffice for 90.5–98.8\% accuracy with cross-fold Jaccard 0.67–0.76.
% (4) Result B — HEADLINE: same-domain donor→target activation patching flips human predictions to AI at 1.07–8.15\%, 10–24× above random-feature baselines, providing the first causal-sufficiency evidence for AI-detection circuits in a frozen encoder.
% (5) Result C: detection circuits are generator-specific (pairwise Jaccard 0.00–0.24) but ~32× above chance; chat-RLHF generators concentrate in layer 12, the only base generator (gpt2) peaks at layer 11.
% (6) Position: complementary to confounder-suppression frameworks (Borile \& Abrate, 2025) — we patch to demonstrate sufficiency rather than ablate to remove bias.
\todo{Write abstract last, after Section 1 is stable.}
\end{abstract}

% =====================================================================
\section{Introduction}
\label{sec:introduction}

% Suggested flow (5–6 paragraphs, ~1 page):
% ¶1 — Hook. Why AI-text detection matters, and why mechanistic understanding (not just detection accuracy) is now the bottleneck. Cite RAID (Dugan et al. 2024), recent detector-failure analyses (Pudasaini et al. 2026).

The rapid development of Large Language Models (LLMs) has made AI-generated text increasingly difficult to distinguish from human writing~\cite{dugan2024raid}. This poses a significant challenge across domains ranging from academic integrity to social media moderation and journalist authenticity. Even though the detection accuracy has improved substantially, with encoder-based classifiers reaching high performance on established benchmarks~\cite{dugan2024raid}, a fundamental question remains unanswered: \textit{What internal representations actually drive these predictions?} Mechanistic interpretability has seen little application to AI-text detection, leaving the field without a causal account of how detection actually works.

% ¶2 — Gap. Existing interpretability for MGT detection is either (a) black-box feature attribution (SHAP/LIME), (b) confounder-removal via neuron suppression in fine-tuned encoders (Borile \& Abrate 2025), or (c) SAE feature analysis on decoders (Kuznetsov et al. 2025). What's missing: a causal-sufficiency test on a frozen encoder.

Existing interpretability work on AI-text detection takes one of three forms: a)~\textit{black-box feature attribution}, which identifies input-level correlates via SHAP or LIME without inspecting model internals~\cite{pudasaini2026why}; b)~\textit{confounder removal}, which suppresses neurons correlated with domain or style artifacts in fine-tuned encoders to improve out-of-distribution accuracy~\cite{borile2025confounding}; or (c)~\textit{sparse autoencoder feature analysis} applied to decoder-based models~\cite{kuznetsov2025sae}. These approaches share a fundamental limitation: they are \textit{observational}. They identify neurons that correlate with detector behavior or remove features that hurt generalization, but none directly tests whether a small set of internal representations is \textit{causally sufficient} to produce AI-text predictions.

% ¶3 — Approach. Frozen \bertbase + Gurnee-style L1-then-L2 sparse probing on RAID across 5 generators (1 base + 4 chat-RLHF). Two interventions on the SAME selected set: mean-ablation (necessity) and same-domain donor→target activation patching (sufficiency).

We address this gap with a causal intervention study on a \textit{frozen} \bertbase. Rather than fine-tuning the model for detection, we train sparse linear probes on pretrained CLS-token activations concatenated across all 12 layers (9{,}216 features total), following the sparse-probing protocol of \citet{gurnee2023haystack}. An L1-regularized logistic regression selects a small stable set of candidate neurons; an independent L2 probe serves as the frozen evaluator for all subsequent experiments. We then apply two causal interventions to the same selected feature set: \textit{mean ablation} which removes these features to test whether they are necessary, and \textit{same-domain donor $\to$ target activation patching} replacing their activations from AI into human samples to test whether they are \textit{sufficient} to flip detector predictions. The patching direction is the key distinction from prior work: rather than removing suspected confounders, we inject candidate detection signal and measure the result.

% ¶4 — Findings preview (mirror abstract): three numbered findings — sparse stable set; causal sufficiency via patching; generator-specific circuits with chat-RLHF vs base layer asymmetry.

We find that it is sufficient. Across five generators spanning one base and four chat-RLHF model families, patching $\sim$50 selected neurons flips human $\to$ AI predictions at 1.07--8.15\%, 10--16$\times$ above a random-feature baseline. This finding provides the first causal-sufficiency evidence for AI-detection circuits in a frozen pretrained encoder. Mean-ablation of the same features causes only 0.03--1.06\,pp accuracy loss, revealing that the detection signal is \textit{redundantly distributed}: easy to inject, hard to erase. Cross-generator analysis further shows that stable feature sets are largely generator-specific (pairwise Jaccard 0.00--0.24) yet cluster $\sim$32$\times$ above chance, with chat-RLHF generators concentrating in layer~12 while the sole base/non-RLHF generator (GPT-2) peaks at layer~11.

% ¶5 — Position vs prior work in one paragraph. Borile \& Abrate differ on 4 axes (frozen vs fine-tuned; CLS vs FFN intermediate; positive-patching vs hard-suppression; RAID vs DAIGT/M4/HC3). Enkhbayar (2025) reports a "causality gap" in observational discriminativity — we close that gap with patching. Kuznetsov et al. 2024 use coordinate-subspace removal, not neuron-level probing.

Our work is most directly related to \citet{borile2025confounding}, who suppress confounding neurons in a fine-tuned encoder to improve out-of-distribution detection, an observational approach that tests necessity-for-confounding rather than causal sufficiency. We discuss the relationship to this and other prior work in Section ~\ref{sec:related}.

% ¶6 — Contributions, bulleted.

\paragraph{Contributions.}
\begin{enumerate}
    \item A sparse-probing characterization of AI-detection signal in frozen \bertbase across five generators on RAID, identifying stable sets of 45-62 neurons ($<$1\% of 9,216 features) sufficient for 90.5-98.8\% detection accuracy with high cross-fold stability (Jaccard 0.67-0.76), with hyperparameters validated via a multi-generator stability sweep.

    \item The first causal-sufficiency test for AI-detection features via same-domain donor $\to$ target activation patching, yielding flip rates 10-16$\times$ above a random-feature baseline across all five generators.

    \item A cross-generator circuit analysis revealing generator-specific stable sets that nonetheless overlap $\sim$32$\times$ above chance, with a chat-RLHF (layer~12) vs.\ base/non-RLHF (layer~11) layer asymmetry.

    \item Leave-one-family-out (LOGO) generalization showing that $\leq$1\% of neurons carries 87-94\% of full-feature detection signal across five unseen model families.
\end{enumerate}

% =====================================================================
\section{Related Work}
\label{sec:related}

% Three subsections, ~2/3 page total. Position carefully — this is where the Borile \& Abrate distinction lives.

\subsection{AI-Generated Text Detection}
% Briefly: supervised classifiers (RoBERTa-based, GPT-2-output detector), zero-shot methods (DetectGPT, Binoculars), benchmark landscape (RAID, M4, HC3). Cite Dugan et al. 2024 for RAID. Keep short — this is not a detection paper.

Supervised detectors fine-tune pretrained language models as binary classifiers on pairs of human and machine-generated text. The GPT-2 output detector~\citep{solaiman2019release} established this paradigm by fine-tuning RoBERTa on GPT-2 samples; subsequent work extended the approach to instruction-tuned and RLHF-aligned generators with larger and more diverse training sets. These classifiers achieve high in-distribution accuracy but degrade substantially on generators or domains unseen during training~\citep{dugan2024raid}.

Zero-shot methods exploit statistical properties of language model distributions without requiring labelled data. DetectGPT~\citep{mitchell2023detectgpt} tests whether candidate text lies at a local maximum of log-probability relative to nearby perturbations, reflecting the tendency of model-generated text to occupy high-probability regions. Binoculars~\citep{hans2024binoculars} uses the cross-perplexity ratio between two models as a calibrated detection score, achieving strong performance across generators. Both approaches require whitebox or API access to a reference model.

Systematic benchmarking has revealed that no single method generalises reliably. RAID dataset~\citep{dugan2024raid} spans 11 generators and 6 domains with adversarial rephrasing attacks; classifiers that perform well in-distribution consistently fail on held-out generator families or domains. This pattern motivates our focus: rather than improving detection accuracy, we ask what internal representations actually drive it.

\subsection{Interpretability of Text Detectors}
% Black-box attribution (Pudasaini et al. 2026 — SHAP-based, finds cross-domain divergence). Confounder-suppression on fine-tuned BERT: Borile \& Abrate (Findings EMNLP 2025) — explicitly differentiate on the 4 axes here. SAE-based feature interp (Kuznetsov et al. 2025, arXiv:2503.03601) — different model class (Gemma-2 decoder). Embedding-subspace removal (Kuznetsov et al. 2024, Findings EMNLP 2024) — geometric, not neuron-level.

The shallowest form of interpretability for AI-text detectors uses input-level attribution. \citet{pudasaini2026why} apply SHAP to explain classifier decisions, finding that the most predictive tokens diverge substantially across domains. This approach characterises the model's decision surface but cannot speak to what internal features underlie it.

More recent work examines representations inside the encoder. \citet{borile2025confounding} identify neurons in a fine-tuned BERT encoder whose activations correlate with domain or style artefacts rather than AI-authorship signal, and suppress them to improve out-of-distribution accuracy. \citet{kuznetsov2024restricted} take a geometric approach, removing a learned embedding subspace to produce a detector more robust to lexical variation. Both works operate on fine-tuned models and test necessity-for-confounding: the question is which features hurt generalisation, not which features are sufficient to drive predictions. Our work differs from \citet{borile2025confounding} on four axes: we use a \textit{frozen} rather than fine-tuned encoder; we probe full-layer concatenated CLS representations rather than intermediate FFN activations; we apply activation patching to test \textit{sufficiency} rather than ablation to identify confounders; and we evaluate on RAID rather than DAIGT/M4/HC3.

\citet{kuznetsov2025sae} apply sparse autoencoders to Gemma-2 decoder activations, extracting semantically interpretable features correlated with AI authorship. This work is complementary but distinct: the model class (autoregressive decoder), representation type (token-position activations rather than CLS pooling), and analysis method (SAE reconstruction) all differ from ours, and the approach remains observational -- features are identified by correlation, not by causal intervention. None of the above provides direct evidence that a small identified set of neurons is \textit{sufficient} to produce AI-text predictions when injected into human-text activations; that is the question we address.

\subsection{Sparse Probing and Activation Patching}
% Sparse probing: Gurnee et al. (TMLR 2023, "Finding Neurons in a Haystack") as canonical reference for L1-then-L2 protocol. Activation patching: Vig et al. 2020 (causal mediation), Meng et al. 2022 (ROME), Heimersheim \& Nanda 2024 (best-practices reference; cite when justifying mean-ablation choice and on-manifold donor patching). Mention Ravindran 2025 (arXiv:2507.09406) as prior use of donor patching with flip rate in safety contexts — but for autoregressive deception, not encoder-based detection.

Sparse probing uses L1-regularised linear classifiers to identify a small set of neurons in a pretrained model that carry task-relevant information. \citet{gurnee2023haystack} introduced and systematically evaluated the L1-then-L2 protocol: an L1 probe selects candidate neurons, and a re-fitted L2 probe on those neurons provides unbiased accuracy estimates. Their study demonstrates that individual neurons in GPT-2 and LLaMA encode grammatical, factual, and linguistic properties, and that the selected sets are stable across seeds and data subsets. We adopt this protocol directly, applying it to \bertbase CLS-token activations concatenated across all 12 layers.

Activation patching replaces internal representations in a target forward pass with those from a donor pass, then measures the effect on model outputs. \citet{vig2020causal} introduced causal mediation analysis for language models, using interventions on attention heads to trace gender bias; \citet{meng2022rome} applied patching at MLP layers in GPT to localise factual associations. \citet{heimersheim2024patching} provide a systematic treatment of patching variants and recommend donor patching -- replacing with activations from a real input rather than a zero or mean vector -- to avoid out-of-distribution artefacts, and mean ablation as a conservative necessity baseline. We follow both recommendations: selected neurons are mean-ablated for the necessity test and replaced with activations from same-domain AI donors for the sufficiency test.

The closest methodological precedent for our flip-rate framing is \citet{ravindran2025adversarial}, who apply donor-style activation patching to test whether deceptive activations from an aligned model can cause unsafe outputs in a target pass. The setting differs substantially -- autoregressive decoder, safety alignment rather than AI-text detection -- but the core intervention and the use of prediction-flip rate as the primary metric are shared. To our knowledge, no prior work applies this methodology to encoder-based AI-text detection.

% =====================================================================
\section{Setup}
\label{sec:setup}

% ~3/4 page. Goal: a reader can reproduce. All numbers from §2 of progress report.

\subsection{Data: RAID}
\label{sec:setup-data}
% RAID benchmark (Dugan et al. 2024): 11 generators × 8 domains. Headline run uses 6 generators chosen for architectural diversity:
% gpt4, gpt2, mpt, mistral-chat, llama-chat, cohere-chat (2 pure-base + 1 SFT-only + 3 SFT+RLHF, spanning 5 model families).
% 6 domains: abstracts, books, news, reddit, reviews, wiki.
% Balanced subsampling: 7,500 samples per generator (3,750 per class), justified in §\ref{sec:setup-stability}.

RAID~\citep{dugan2024raid} provides human-written and machine-generated texts across 11 generators and 8 domains. We use six of the eight domains: abstracts, books, news, Reddit posts, reviews, and Wikipedia. We select six generators for architectural diversity: GPT-4 (\texttt{gpt4}), GPT-2 (\texttt{gpt2}), MPT-7B (\texttt{mpt}), Mistral-7B-Instruct (\texttt{mistral-chat}), Llama-2-Chat (\texttt{llama-chat}), and Cohere Command (\texttt{cohere-chat}), spanning two pure-base generators (gpt2, mpt), one SFT-only generator (mistral-chat), and three SFT+RLHF/preference-trained generators (gpt4, llama-chat, cohere-chat), across five model families. We subsample 7{,}500 examples per generator (3{,}750 human, 3{,}750 AI) balanced across the six RAID domains, yielding 45{,}000 samples in total.

\subsection{Model and Activations}
\label{sec:setup-model}
% Frozen \bertbase (12 layers, hidden dim 768). No fine-tuning anywhere.
% CLS-token activation extracted at every layer; concatenation gives a 12 × 768 = 9,216-dim feature vector per sample.
% Note: we deliberately use the frozen pretrained model, NOT a detection-fine-tuned one — this isolates representation present in pretraining rather than learned during detection training. (Contrast with Borile \& Abrate.)

We use \bertbase~\citep{devlin2019bert} with all weights frozen throughout. For each input text we extract the CLS-token hidden state at every transformer layer $l \in \{1, \ldots, 12\}$, giving a per-layer vector $\mathbf{h}^{(l)} \in \mathbb{R}^{768}$. The full feature representation is the concatenation
\[
  \mathbf{z} = \bigl[\mathbf{h}^{(1)};\, \ldots;\, \mathbf{h}^{(12)}\bigr] \in \mathbb{R}^{9{,}216}.
\]
We refer to individual dimensions of $\mathbf{z}$ as \textit{neurons}; neuron index $i$ resides in layer $\lfloor i / 768 \rfloor + 1$ and position $i \bmod 768$ within that layer. Using the frozen pretrained model (rather than a fine-tuned detector) ensures that any informative features we identify are present in general-purpose BERT representations, not artefacts of task-specific fine-tuning.

\subsection{Protocol: L1 Selection, L2 Evaluation}
\label{sec:setup-protocol}
% Per (seed, fold):
% (1) L1 logistic regression at fixed C=0.005 on the train fold (liblinear solver, max_iter=1000) → SelectionResult: ~$\nselected$ neurons (45–92 depending on generator/cell), ranking by |coef|.
% (2) Independent L2 logistic regression on ALL 9,216 features (lbfgs) → eval probe. The eval probe is FIXED across evaluators.
% (3) Each evaluator (ablation/patching/AUC-comparison/restricted) sees test-fold activations + frozen selection + frozen eval probe.
% Cross-validation: StratifiedKFold on (label × domain), 5 folds × 3 seeds [42, 123, 456] = 15 cells per (generator, experiment).

Following \citet{gurnee2023haystack}, for each (seed, fold) cell we fit two independent probes on the training split. (1)~\textit{L1 selection probe}: L1-regularised logistic regression ($C = 0.005$, liblinear solver) on all 9{,}216 features. Neurons with non-zero coefficients form the \textit{candidate set} $\mathcal{S}$; neurons appearing in $\mathcal{S}$ in at least 50\% of the 15 cells are \textit{stable neurons}. (2)~\textit{L2 evaluation probe}: L2-regularised logistic regression (lbfgs) on the full $\mathbf{z}$. This probe is held frozen for all downstream experiments to ensure that intervention effects are measured against a fixed decision function.

We use 5-fold stratified cross-validation stratified by (label $\times$ domain) with seeds $\{42, 123, 456\}$, yielding 15 cells per (generator, experiment). All reported metrics are 15-cell means $\pm$ standard deviation unless otherwise noted.

\subsection{Stability-Driven Hyperparameter Selection}
\label{sec:setup-stability}
% Brief: K=15 stratified subsamples × 4 C values × 7 N values × 3 generators (gpt4, llama-chat, mistral). Pairwise Jaccard at C=0.005, N=7,500 is 0.72–0.75 with min-over-3-gen 0.724 — most-stable C value at every N for every generator. Full sweep deferred to Appendix \ref{app:stability}.

The regularisation strength $C = 0.005$ and sample size $N = 7{,}500$ were selected via a multi-generator stability sweep conducted prior to any results analysis, following the criterion of \citet{gurnee2023haystack}. We drew $K = 15$ stratified subsamples from each of three generators (\texttt{gpt4}, \texttt{llama-chat}, \texttt{mistral-chat}), computed mean pairwise Jaccard similarity between selection sets across subsamples, and chose the $(C, N)$ pair that maximised the minimum Jaccard across generators. At $C = 0.005$, $N = 7{,}500$, the minimum-over-generators Jaccard is 0.724 --- the highest achieved by any $C$ at every $N$ tested. The full sweep is reported in Appendix~\ref{app:stability}.

% =====================================================================
\section{Sparse Detection Circuits in Frozen BERT}
\label{sec:results-sparse}

% ~1/2 page. The "small stable set suffices" result.
% Per-cell accuracy/F1/AUC table → appendix.

Table~\ref{tab:sparse-summary} reports sparse-probe results across all six generators. Between 45 and 62 neurons --- 0.5--0.9\% of the 9{,}216-feature CLS representation --- are stable across all 15 cells, yet the L2 evaluation probe achieves 90.5--98.8\% validation accuracy (AUC-ROC 0.968--0.999). Selection is highly reproducible: pairwise fold Jaccard ranges 0.67--0.76 across 105 fold-pairs per generator, compared to a random-null expectation of $\sim$0.003--0.005 for sets of this size drawn from 9{,}216 features, placing within-generator stability at 130--200$\times$ above chance. The selection pool is narrow: between 91 neurons ever selected (llama-chat) and 146 (gpt2), indicating that the probe consistently returns to a small region of the feature space.

Detection difficulty varies systematically across generators. GPT-4 is easiest to detect (98.8\% val acc, AUC 0.999), followed by LLaMA-2-Chat (98.6\%) and Mistral-7B-Instruct (97.1\%); all three are instruction-tuned and produce outputs that diverge markedly from human writing in BERT's frozen representation space. MPT-7B (97.2\%) is close behind despite being a pure-base generator, with 52 stable neurons forming a compact and highly reproducible selection set. GPT-2 sits at the mid-point (95.9\%) with the loosest fold Jaccard (0.670) and largest observed pool (146 neurons) among all six generators, suggesting a noisier but stable feature landscape for this base model. Cohere-Command-Chat is the structural difficulty floor at 90.5\% val accuracy and AUC 0.968, with the highest standard deviation on $\nselected$ (CV 4.3\%), consistent with its smaller per-axis signal magnitude documented in Section~\ref{sec:results-crossgen}.

\begin{table}[t]
\centering
\small
\begin{tabular}{lrrrr}
\toprule
Generator & Val acc & $\nselected$ & Jaccard & $n_{\text{stable}}$ / $n_{\text{obs}}$ \\
\midrule
gpt4         & $0.988 \pm 0.003$ & $74.1 \pm 2.9$ & $0.763 \pm 0.037$ & 60 / 106 \\
gpt2         & $0.959 \pm 0.003$ & $84.7 \pm 4.1$ & $0.670 \pm 0.042$ & 60 / 146 \\
mpt          & $0.972 \pm 0.003$ & $61.3 \pm 3.1$ & $0.753 \pm 0.040$ & 52 / 93  \\
mistral-chat & $0.971 \pm 0.004$ & $57.8 \pm 2.9$ & $0.724 \pm 0.046$ & 45 / 99  \\
llama-chat   & $0.986 \pm 0.003$ & $63.6 \pm 1.7$ & $0.756 \pm 0.040$ & 48 / 91  \\
cohere-chat  & $0.905 \pm 0.008$ & $72.3 \pm 3.1$ & $0.749 \pm 0.039$ & 62 / 113 \\
\bottomrule
\end{tabular}
\caption{Sparse-probe results across six generators: 15-cell means $\pm$ std (5 folds $\times$ 3 seeds). Val acc is the L2 evaluation probe accuracy (all 9{,}216 features). $\nselected$: mean candidate set size per cell. Jaccard: pairwise fold Jaccard over 105 fold-pairs (mpt estimated analytically from per-neuron stability frequencies). $n_{\text{stable}}$: neurons in $\geq$50\% of cells; $n_{\text{obs}}$: total distinct neurons selected at least once.}
\label{tab:sparse-summary}
\end{table}

% =====================================================================
\section{Causal Sufficiency via Activation Patching}
\label{sec:results-patching}

% ~1 page. THE HEADLINE SECTION.
% Use this section to explicitly distinguish necessity (ablation, §\ref{sec:results-necessity}) from sufficiency (this section).

This section tests whether the selected neurons are \textit{causally sufficient} to produce AI-text predictions. If transplanting their activations from a real AI sample into a human sample causes the evaluation probe to flip its prediction, the selected set carries the detection signal rather than merely correlating with it. This is a strictly stronger claim than probe accuracy alone.

\subsection{Patching Protocol}
\label{sec:patching-protocol}
% For each fold, restrict to human samples (label=0) that have at least one same-domain AI sample in the test fold ("valid_humans").
% For each of n_shuffles=20 random within-domain donor assignments:
%   patched_acts = base_acts; patched_acts[:, top_k_indices] = donor_acts[:, top_k_indices]
%   flip = (eval_probe.predict(patched) == 1) AND (eval_probe.predict(base) == 0)
% Flip rate = flips / |valid_humans|. Mean/std over 20 shuffles, then over 15 cells.
% Random-k baseline: n_random_draws=20 random k-subsets, each run with 20 shuffles.
% Note: same-domain donor sampling controls for domain shift (Heimersheim & Nanda 2024 on-manifold guidance).

For each (seed, fold) cell we identify \textit{valid human} samples: test-fold human texts that share a domain with at least one test-fold AI sample. For each of $n_{\text{shuffles}} = 20$ random within-domain donor permutations we construct a patched activation $\tilde{\mathbf{z}}$ by replacing the selected neurons with those from the AI donor:
\[
  \tilde{\mathbf{z}}_i =
  \begin{cases}
    \mathbf{z}^{\text{donor}}_i & i \in \mathcal{S} \\
    \mathbf{z}^{\text{human}}_i & i \notin \mathcal{S}
  \end{cases}
\]
A \textit{flip} occurs when the frozen evaluation probe classifies $\tilde{\mathbf{z}}$ as AI but the unpatched $\mathbf{z}^{\text{human}}$ as human. The flip rate is the fraction of valid human samples that flip, averaged over the 20 donor permutations and then over the 15 cells. Donor sampling is restricted to the same domain as the target, following the on-manifold recommendation of \citet{heimersheim2024patching} to avoid injecting out-of-distribution activations.

The \textit{random-$k$ baseline} draws 20 random subsets of size $k$ from the 9{,}216 features and runs the identical procedure, providing a null that controls for the number of patched dimensions. We sweep $k \in \{1, 5, 10, 20, 50, \text{full}\}$ where $k = \text{full}$ equals the cell's actual $|\mathcal{S}|$ (53--92 neurons depending on generator and cell).

\subsection{Flip-Rate Results}
\label{sec:patching-results}
% Table: k-sweep (k = 1, 5, 10, 20, 50, full) × 5 generators, selected flip rate only.
% Headline at k=full: gpt4 1.23\%, gpt2 3.13\%, mistral 1.73\%, llama 1.07\%, cohere 8.15\% — monotone in k for every generator.
% Ratio vs random-k baseline at k=full: 10× (gpt2), 12× (llama), 11× (cohere), 23× (mistral), ~125× (gpt4 — caveat: random baseline ~10^-4 inflates the ratio with high variance; report selected-flip column as the robust signal).
% Per-domain breakdown table → appendix or compact figure.
% Interpretation: "selected neurons are sufficient to push human predictions toward AI on all five generators." Cohere is most easily perturbed (probe least confident, 90.5\% baseline); llama least (1.07\%).

Table~\ref{tab:patching-sweep} reports selected flip rates across the $k$-sweep. The result is consistent across all six generators: flip rate is monotonically increasing in $k$, and the selected set outperforms the random-$k$ baseline at every $k$ value. At $k = \text{full}$, patching the selected neurons causes 1.07--8.15\% of human samples to be reclassified as AI-generated. Against a random baseline of 0.08--0.78\% at the same $k$, this yields a ratio of 10--23$\times$ on five of six generators. For gpt4 the ratio is $\sim$125$\times$ but with standard deviation $\approx$265, because the random-$k$ flip rate is near zero ($\sim 8 \times 10^{-4}$) and small denominator fluctuations dominate; the selected flip rate of 1.23\% is the robust signal there. MPT-7B falls at the lower end of the range (1.59\%, $\sim$10$\times$ above random), consistent with its moderate detection difficulty.

Cohere-Command-Chat is the most patchable generator (8.15\% at $k = \text{full}$, $\pm$1.32\%), consistent with its lowest baseline accuracy (90.5\%): a probe operating nearer its decision boundary crosses it under a smaller perturbation. LLaMA-2-Chat is the hardest to flip (1.07\%, $\pm$0.38\%) despite its high baseline accuracy (98.6\%), suggesting its detection signal is more broadly distributed. Flip rate scales approximately log-linearly with $k$ across all generators, indicating that each added neuron contributes incrementally rather than the effect concentrating in a few critical units.

\begin{table}[t]
\centering
\small
\begin{tabular}{lrrrrrr}
\toprule
$k$ & gpt4 & gpt2 & mpt & mistral & llama & cohere \\
\midrule
1    & 0.02 & 0.05 & 0.10 & 0.06 & 0.07 & 0.13 \\
5    & 0.09 & 0.17 & 0.25 & 0.20 & 0.14 & 0.58 \\
10   & 0.15 & 0.29 & 0.43 & 0.34 & 0.27 & 0.99 \\
20   & 0.35 & 0.58 & 0.79 & 0.56 & 0.49 & 2.02 \\
50   & 0.80 & 1.71 & 1.52 & 1.48 & 0.86 & 5.43 \\
\midrule
full & $1.23_{\pm 0.34}$ & $3.13_{\pm 0.56}$ & $1.59_{\pm 0.23}$ & $1.73_{\pm 0.28}$ & $1.07_{\pm 0.38}$ & $8.15_{\pm 1.32}$ \\
\bottomrule
\end{tabular}
\caption{Selected flip rate (\%) across the $k$-sweep, 15-cell means (5 folds $\times$ 3 seeds). $k = \text{full}$ equals each cell's actual $|\mathcal{S}|$ (53--92 neurons); std shown only for $k = \text{full}$. mpt $k < \text{full}$ values from a representative fold. Flip rate is monotonically increasing in $k$ for every generator.}
\label{tab:patching-sweep}
\end{table}

\subsection{Robustness Checks}
\label{sec:patching-robustness}
% (a) Random-donor (cross-domain) baseline: report flip-rate gap vs same-domain donor → demonstrates the effect is not driven by domain transfer.
% (b) Zero-ablation comparison alongside mean-ablation in §\ref{sec:results-necessity} — Heimersheim \& Nanda 2024 on off-manifold concerns for mean.
\todo{Run cross-domain donor baseline before submission: replace same-domain donors with random cross-domain AI donors and measure flip rate gap vs same-domain. Expected to confirm the effect is not driven by domain-level activation shifts. Also consider reporting zero-ablation alongside mean-ablation in \S\ref{sec:results-necessity} per \citet{heimersheim2024patching} off-manifold concern.}

% =====================================================================
\section{Necessity via Mean Ablation}
\label{sec:results-necessity}

% ~1/2 page. Honest framing — "not individually necessary on 4/5 generators."
% Mean-ablation of top-k selected neurons (replace with train-fold mean to avoid leakage); score the SAME L2 eval probe on test.
% k-sweep table reused from §3.2 of progress report.
% Headline:
% — 4/5 generators: ablation drop at k=full sits at 0.03–0.11pp, within noise of the 0.06pp random-k baseline;
% — cohere-chat exception: 1.06 ± 0.54pp drop, ~10× random baseline, monotone from k=5.
% Interpretation tied to §\ref{sec:patching-results}:
% Ablation tests necessity against a probe with all 9,216 features — removing ≤1\% of them leaves redundant correlates.
% Patching tests sufficiency by injecting wrong-class signal — competes with the bulk of the network.
% No contradiction: a redundantly distributed representation is exactly what gives "small drop on ablation, large flip on patching."
% Methodology caveat paragraph (1–2 sentences): mean-ablation is off-manifold (Heimersheim \& Nanda 2024). We additionally report zero-ablation in §\ref{sec:patching-robustness} and treat the necessity result as conservative.

To complement the sufficiency test we ask the converse question: is the selected set \textit{necessary}? We answer this via mean ablation, replacing the selected neurons in the test-fold activations with their training-fold mean (computed on the training split to avoid leakage) and re-evaluating the frozen L2 probe. Table~\ref{tab:ablation-full} reports accuracy drop at $k = \text{full}$ alongside the random-$k$ baseline.

For four of five generators the result is negative. At $k = \text{full}$, ablating the selected set causes accuracy to drop by only 0.03--0.11\,pp --- within noise of the 0.00--0.06\,pp random-$k$ baseline. The selected neurons are not individually necessary: the L2 evaluation probe, which uses all 9{,}216 features, retains thousands of redundant correlates that compensate for the removed subset.

Cohere-Command-Chat is the exception. Its full-set ablation drop is 1.06\,$\pm$\,0.54\,pp --- roughly 18$\times$ the random baseline of 0.06\,pp --- and the drop grows monotonically from $k = 5$ upward, unlike the other four generators where the curve is flat. This pattern is consistent with the other cohere signals (lowest baseline accuracy, smallest mean-difference norm, largest probe weight norm) and suggests its detection representation is \textit{partially} localised rather than fully distributed.

There is no contradiction between the small ablation drops here and the large patching flip rates in Section~\ref{sec:results-patching}. Ablation tests necessity \textit{against a probe that has access to all other features}: removing $\leq$1\% of 9{,}216 features leaves thousands of correlated dimensions intact. Patching tests sufficiency by \textit{injecting wrong-class signal that competes with the remainder of the network}. A representation that is redundantly distributed is exactly what would produce small ablation drops alongside meaningful patching effects. We interpret the joint result as: the selected neurons carry genuine AI-detection signal, but the same signal is also distributed across the broader feature space.

A methodological note: mean ablation replaces activations with out-of-distribution values that fall off the natural data manifold, which can introduce artefacts \citep{heimersheim2024patching}. We treat the necessity result as conservative and plan to report zero-ablation as a robustness check (Section~\ref{sec:patching-robustness}).

\begin{table}[t]
\centering
\small
\begin{tabular}{lrrr}
\toprule
Generator & Baseline acc & Selected drop ($k{=}\text{full}$) & Random drop ($k{=}\text{full}$) \\
\midrule
gpt4         & 0.988 & $+0.03 \pm 0.12$\,pp & $+0.01$\,pp \\
gpt2         & 0.959 & $+0.11 \pm 0.31$\,pp & $+0.00$\,pp \\
mpt          & 0.972 & $+0.05 \pm 0.11$\,pp & $+0.00$\,pp \\
mistral-chat & 0.971 & $+0.10 \pm 0.21$\,pp & $+0.00$\,pp \\
llama-chat   & 0.986 & $+0.11 \pm 0.14$\,pp & $+0.02$\,pp \\
cohere-chat  & 0.905 & $\mathbf{+1.06 \pm 0.54}$\,pp & $+0.06$\,pp \\
\bottomrule
\end{tabular}
\caption{Mean-ablation accuracy drop at $k = \text{full}$, 15-cell means $\pm$ std. Drop = baseline $-$ ablated accuracy (positive = accuracy decreases). Random drop is the mean over 20 random $k$-subsets of the same size. Cohere-chat is the only generator with a drop substantially above the random baseline.}
\label{tab:ablation-full}
\end{table}

% =====================================================================
\section{Cross-Generator Circuit Geometry}
\label{sec:results-crossgen}

% ~3/4 page. The "circuits are generator-specific but overlap above chance" + "RLHF/base layer asymmetry" finding.

\subsection{Layer Distribution of Stable Neurons}
\label{sec:crossgen-layers}
% Stack-bar or table: 12 layers × 5 generators, count of stable neurons.
% Highlight: chat-RLHF generators concentrate 30–36\% of stable neurons in layer 12.
% gpt2 (only base / non-RLHF generator): peaks at layer 11 (18.3\%); only 13.3\% in layer 12.
Table~\ref{tab:layer-dist} shows the layer-by-layer distribution of stable neurons across all six generators. A clear asymmetry emerges between the four instruction-tuned generators and the two pure-base generators. For gpt4, mistral-chat, llama-chat, and cohere-chat, the single most populated layer is layer~12, which accounts for 30--36\% of each generator's stable set (16--19 neurons). Both pure-base generators are exceptions: GPT-2 peaks at layer~11 with 11 stable neurons (18.3\% of its stable set) and has only 13.3\% in layer~12; MPT-7B has just 2 stable neurons in layer~12 (3.8\%) and distributes its signal more broadly across earlier layers (layers 1, 8, and 9 dominate).

Across both pure-base generators, the shared pattern is markedly low layer-12 concentration (3.8--13.3\%) compared to any instruction-tuned generator (30--36\%). This 2-vs-4 contrast suggests that BERT's final layer functions as a generator-specific classification register preferentially activated by instruction-tuned outputs, while base-model outputs engage BERT's earlier and mid-depth representations more strongly. The specific peak layer differs between gpt2 (L11) and mpt (L1/L8), indicating the effect is not a rigid ``base generators peak at L11'' rule but rather a consistent \textit{avoidance} of L12 by non-instruction-tuned generators.

\begin{table}[t]
\centering
\small
\begin{tabular}{lrrrrrr}
\toprule
Layer & gpt4 & gpt2 & mpt & mistral & llama & cohere \\
\midrule
1--10  & 36 & 41 & \textbf{44} & 24 & 30 & 40 \\
11     &  6 & \textbf{11} &  6 &  5 &  2 &  3 \\
12     & \textbf{18} &  8 &  2 & \textbf{16} & \textbf{16} & \textbf{19} \\
\midrule
\% in L12 & 30.0 & 13.3 & 3.8 & 35.6 & 33.3 & 30.6 \\
\bottomrule
\end{tabular}
\caption{Stable-neuron counts by layer group across six generators. Bold indicates the peak layer group for each generator. Instruction-tuned generators (gpt4, mistral, llama, cohere) concentrate 30--36\% of stable neurons in layer~12; both pure-base generators (gpt2, mpt) have $\leq$14\% in layer~12.}
\label{tab:layer-dist}
\end{table}

\subsection{Stable-Set Overlap Across Generators}
\label{sec:crossgen-jaccard}
% Pairwise Jaccard matrix (10 pairs across 5 generators).
% Absolute Jaccard: 0.00–0.24.
% Random-null Jaccard (analytic, two random size-K_A, K_B sets in 9,216 features): ~0.003.
% → Mean over 10 pairs is ~32× chance.
% Strongest pair: cohere-mistral (J=0.244, 86× chance). Weakest: gpt2-gpt4 (J=0.000, 0× chance).
% Interpretation: detection circuits are generator-specific in absolute terms but overlap far above chance baseline — supports a partial shared substrate.

Table~\ref{tab:jaccard} reports pairwise Jaccard similarity between stable sets across all fifteen generator pairs. In absolute terms, overlap is low: Jaccard ranges from 0.00 to 0.24, consistent with generator-specific circuits. However, the appropriate baseline is not zero but the expected Jaccard for two random sets of the same sizes drawn from 9,216 features, which is approximately 0.003 per pair. Relative to this null, the mean observed Jaccard of 0.073 is $\sim$24$\times$ above chance, indicating a partial shared substrate despite the generator-specific structure.

The data reveal a clear bipartite structure. Among instruction-tuned generators, pairwise Jaccard ranges 0.08--0.24 (27--86$\times$ chance), with the strongest pairs being cohere--mistral (86$\times$) and llama--mistral (64$\times$). Among the two pure-base generators, mpt and gpt2 share 12 neurons (J\,=\,0.120, 40$\times$ chance) --- higher than several instruction--instruction pairs --- confirming that the base models form a coherent sub-cluster. In contrast, every base-vs-instruction-tuned pair has near-zero or zero overlap: mpt has J\,=\,0.000 with both gpt4 and cohere-chat, and J\,$\leq$\,0.011 with mistral and llama. Three pairs now have zero overlap (mpt--gpt4, mpt--cohere, gpt2--gpt4), all involving at least one base generator. This base-vs-instruction-tuned partition explains the bulk of the cross-generator circuit geometry.

\begin{table}[t]
\centering
\small
\begin{tabular}{lrr}
\toprule
Pair & Jaccard & $\times$ chance \\
\midrule
cohere -- mistral & 0.244 & 86$\times$ \\
llama -- mistral  & 0.163 & 64$\times$ \\
gpt4 -- mistral   & 0.141 & 51$\times$ \\
gpt4 -- llama     & 0.113 & 39$\times$ \\
\textbf{mpt -- gpt2}       & \textbf{0.120} & \textbf{40$\times$} \\
cohere -- gpt4    & 0.109 & 33$\times$ \\
cohere -- llama   & 0.078 & 27$\times$ \\
gpt2 -- mistral   & 0.050 & 18$\times$ \\
cohere -- gpt2    & 0.034 & 10$\times$ \\
gpt2 -- llama     & 0.019 &  7$\times$ \\
mpt -- mistral    & 0.010 &  4$\times$ \\
mpt -- llama      & 0.010 &  4$\times$ \\
gpt2 -- gpt4      & 0.000 &  0$\times$ \\
mpt -- gpt4       & 0.000 &  0$\times$ \\
mpt -- cohere     & 0.000 &  0$\times$ \\
\midrule
Mean              & 0.073 & $\sim$24$\times$ \\
\bottomrule
\end{tabular}
\caption{Pairwise Jaccard similarity between stable sets (15 generator pairs). $\times$\,chance is observed Jaccard divided by the analytic random-null Jaccard for two sets of the same sizes drawn from 9,216 features ($\approx$0.003 per pair). The mpt--gpt2 pair (bold) is the highest base$\times$base overlap; all base$\times$instruction-tuned pairs cluster at the bottom.}
\label{tab:jaccard}
\end{table}

\subsection{CAV Diagnostics}
\label{sec:crossgen-cav}
% Single-fit per-generator on stable sets: cosine sim between mean-difference CAV and L1 weight, LR accuracy on CAV, ‖mean diff‖, ‖LR weight‖.
% Cohere has lowest CAV cosine (0.057), smallest mean-diff norm (4.69), largest LR weight norm (107) — consistent with "weaker per-axis signal, compensated by larger weights."
% Brief paragraph; full table can go inline.

To characterise the geometry of the detection signal within each stable set we compute a concept activation vector (CAV) as the mean difference between AI and human activations restricted to the stable neurons, and compare it to the L1 probe weight vector. Table~\ref{tab:cav} reports cosine similarity between the CAV and the probe weight, logistic-regression accuracy on the CAV projection, and the norms of both vectors.

Cohere-chat has the lowest CAV cosine similarity (0.057), the lowest CAV-LR accuracy (0.903), and by far the largest probe weight norm (107), consistent with a diffuse signal compensated by large weights. MPT-7B presents an interesting contrast: it has the lowest cosine (0.065, near cohere's) but the \textit{largest} mean-difference norm (9.52) and smallest weight norm (53.9) among all six generators --- suggesting that mpt's detection signal is strong in per-axis magnitude but geometrically misaligned with the probe direction relative to instruction-tuned generators. GPT-2 also shows a low CAV cosine (0.071) and elevated weight norm (69), placing both base generators at the low-cosine end. Instruction-tuned generators (gpt4, llama-chat, mistral-chat) cluster at higher cosine (0.11--0.14) with moderate weight norms (47--59). CAV diagnostics are currently single-fit per generator and carry no standard errors; wrapping them in the multi-seed loop is a planned runner enhancement.

\begin{table}[t]
\centering
\small
\begin{tabular}{lrrrr}
\toprule
Generator & Cosine (CAV, $w$) & LR acc & $\|\Delta\mu\|$ & $\|w\|$ \\
\midrule
gpt4         & 0.140 & 0.985 &  7.34 &  47.9 \\
gpt2         & 0.071 & 0.969 &  6.84 &  69.1 \\
mpt          & 0.065 & 0.971 &  \textbf{9.52} &  \textbf{53.9} \\
mistral-chat & 0.111 & 0.963 &  6.38 &  58.8 \\
llama-chat   & 0.129 & 0.990 &  6.57 &  47.2 \\
cohere-chat  & \textbf{0.057} & \textbf{0.903} &  4.69 & 107.0 \\
\bottomrule
\end{tabular}
\caption{CAV diagnostics on per-generator stable sets. Cosine: similarity between the mean-difference CAV and the L1 probe weight vector. LR acc: logistic-regression accuracy using only the CAV projection. $\|\Delta\mu\|$: norm of the AI$-$human mean-difference vector (bold = largest). $\|w\|$: L1 probe weight norm (bold = smallest). Single-fit per generator; no standard errors available.}
\label{tab:cav}
\end{table}

% =====================================================================
\section{Discussion}
\label{sec:discussion}

% ~1/2 page. Three threads to land.

% Thread 1 — Necessity vs sufficiency reconciled.
% Redundancy is the natural reading: detection signal is encoded by a small stable set sufficient to flip predictions but distributively backed up across the rest of the residual stream. Ablation under-states the role of selected features because the L2 probe uses all 9,216 dimensions; patching reveals the role because it injects wrong-class signal that must compete with the redundancy.

% Thread 2 — The cohere and gpt2 outliers as evidence, not noise.
% Cohere: most easily flipped (8.15\%), largest LR weight norm, smallest mean-difference norm, clearest necessity contribution. Reading: dispersed signal with weaker per-axis encoding.
% gpt2: zero overlap with gpt4 stable set, layer-11 peak, second-largest restricted-probe gap. Reading: base/non-RLHF generators encode detection signal earlier and in a more dispersed register than chat-RLHF generators.
% Both anomalies are evidence FOR generator-specific representational geometry, not against the central claim.

% Thread 3 — Position vs Borile \& Abrate (2025), one paragraph.
% Their work removes confounder neurons in fine-tuned BERT to improve OOD accuracy; ours patches detection neurons in frozen BERT to demonstrate causal sufficiency. The two are complementary: their setting asks "what should we remove for a better detector?" Ours asks "what does the pretrained model already encode that supports detection?"

\paragraph{Redundancy and sufficiency are not in tension.}
The joint pattern across Sections~\ref{sec:results-patching} and~\ref{sec:results-necessity} --- large patching flip rates alongside negligible ablation drops on 4/5 generators --- has a natural reading: AI-detection signal is \textit{redundantly distributed}. The selected neurons carry genuine causal signal (evidenced by patching), but the same information is backed up across the remaining 9,000+ features in the residual stream (evidenced by the small ablation drop). Ablation understates the role of the selected set because removing $\leq$1\% of a 9,216-dimensional representation leaves thousands of correlated features intact. Patching reveals that role precisely because it injects wrong-class signal that must compete with the rest of the network --- a harder test. The phrase ``easy to inject, hard to erase'' captures this geometry: the selected neurons are sufficient but not necessary.

\paragraph{Outliers as structural evidence.}
Cohere-Command-Chat deviates from other generators on multiple axes: weakest per-axis signal ($\|\Delta\mu\| = 4.69$, CAV cosine 0.057), largest probe weights (norm 107), most patchable output (8.15\% flip rate), and the only generator with a substantial ablation cost (1.06\,pp). Together these point to a signal encoded diffusely with small per-feature magnitude --- the probe compensates via large weights, patching is easy because the decision boundary is near, and ablation matters more because fewer redundant correlates exist. The two pure-base generators, GPT-2 and MPT-7B, form a coherent structural cluster: both have near-zero or zero stable-set overlap with instruction-tuned generators (J\,$\leq$\,0.05), yet they share 12 neurons with each other (J\,=\,0.120, 40$\times$ chance); both avoid layer~12 (3.8\% and 13.3\% respectively, versus 30--36\% for instruction-tuned models); and both exhibit the two largest restricted-probe gaps of the six generators (mpt 4.68\,pp, gpt2 3.63\,pp, versus $\leq$1.1\,pp for the others). This consistency across two independent base models strengthens the interpretation that post-training alignment leaves a layer-12 footprint that is absent in base-only decoding.

\paragraph{Relation to confounder-suppression approaches.}
\citet{borile2025confounding} identify and suppress neurons correlated with domain or style artefacts in a fine-tuned BERT encoder to improve out-of-distribution detection. Our work is complementary along the intervention axis: their question is ``which neurons should we remove to make the detector more robust?'', while ours is ``which neurons does the pretrained model already use, and are they causally sufficient?'' The two approaches converge on the same empirical substrate --- a small subset of BERT neurons carries disproportionate detection-relevant information --- but reach it from opposite directions, necessity-for-robustness versus sufficiency-for-detection.

% =====================================================================
\section{Limitations}
\label{sec:limitations}

% Required by EMNLP. Honest paragraph (or short list). Do NOT bury anything that survived to the discussion.
% Items to cover:
% — Frozen \bertbase only; results may differ on encoder-decoder or larger encoders (RoBERTa-large, DeBERTa).
%   \todo{If RoBERTa replication added (Tier-2), tone this down accordingly.}
% — Mean-ablation is off-manifold (Heimersheim \& Nanda 2024); we report zero-ablation as a check (§\ref{sec:patching-robustness}) but a full resample-ablation comparison is future work.
% — Cross-domain generalization: \todo{If LODO is run, summarize result here in one sentence; if not, state explicitly that within-distribution analysis cannot rule out domain-correlated artifacts and remove any cross-domain claims from the abstract.}
% — RAID adversarial subsets are not evaluated; robustness of the discovered circuits to paraphrase/perturbation attacks is left to future work.
% — N=2 base/non-RLHF generators (gpt2, mpt) for the layer-asymmetry claim. Both avoid L12 (3.8% and 13.3% vs 30-36% for instruction-tuned). Peak layer differs (gpt2: L11, mpt: L1/L8), so claim is framed as L12-avoidance, not L11-peaking.
% — Single encoder family (BERT); SAE-based feature analysis (cf. Kuznetsov et al. 2025 on Gemma-2) is a natural extension.

\paragraph{Single encoder architecture.}
All experiments use frozen \bertbase. Whether the stable-neuron structure and patching effects extend to larger encoders (RoBERTa-large, DeBERTa-v3) or encoder-decoder architectures is an open question. The frozen setting is a deliberate design choice --- it isolates what the pretrained representation already encodes --- but it means results cannot be directly compared to fine-tuned detectors without controlling for this variable.

\paragraph{Mean ablation is off-manifold.}
Replacing selected neurons with their training-mean pushes activations off the natural data manifold, which can introduce artefacts unrelated to the detection signal~\citep{heimersheim2024patching}. We treat the necessity results as conservative lower bounds and plan to report zero-ablation as a robustness check (Section~\ref{sec:patching-robustness}); a full resample-ablation comparison remains future work.

\paragraph{Layer-asymmetry claim across two base generators.}
The observation that instruction-tuned generators concentrate 30--36\% of stable neurons in layer~12 while pure-base generators have $\leq$14\% is now supported by two independent base models (GPT-2 and MPT-7B). However, the two base generators differ in their peak layer (GPT-2: L11; MPT-7B: L1/L8), so the claim is best stated as L12-avoidance rather than a strict L11-peaking rule. We frame this as a consistent L12 footprint of post-training alignment rather than a mechanistically precise claim about which specific earlier layer is activated.

\paragraph{Adversarial robustness not evaluated.}
RAID includes eleven adversarial rephrasing attacks (paraphrase, homoglyph, synonym substitution, etc.) that substantially reduce detector accuracy~\citep{dugan2024raid}. Whether the stable neurons identified here remain causal under these attacks is left to future work; the present analysis covers only the non-adversarial RAID subset.

\paragraph{Scope of causal claim.}
Activation patching demonstrates that the selected neurons are sufficient to shift probe predictions, but does not constitute a full circuit-level proof of causality in the sense of~\citet{meng2022rome}. The patching intervention is at the level of individual neuron activations, not a complete mechanistic decomposition of the computation. We make no claims about what linguistic properties these neurons encode or why they are informative.

% =====================================================================
\section*{Ethics Statement}
\label{sec:ethics}
% Short paragraph. AI-text detection has dual-use considerations: detectors enable platform integrity but also can be used to suppress legitimate AI-assisted writing. Our work is mechanistic-interpretability oriented and does not propose deployment of a new detector. RAID is publicly released under its own license; no human-subjects data was collected.

% =====================================================================
\section*{Acknowledgments}
\label{sec:ack}
% Anonymized for review. Add at camera-ready.

% =====================================================================
\bibliography{refs}

% =====================================================================
\appendix

\section{Pipeline Details}
\label{app:pipeline}
% Full pipeline diagram from progress report §2.1; runner contract pseudocode; SelectionResult schema; per-experiment YAML configuration summary.

\section{Stability Sweep}
\label{app:stability}
% Full Jaccard × C × N grid from progress report §4.5. Justification of C=0.005 and N=7,500.

\section{Per-Cell Metrics}
\label{app:per-cell}
% Per-(seed, fold) accuracy / F1 / AUC tables for sparse probe and the dependent experiments. The headline run identifier is ``20260426\_20\_full\_c0005\_n7500\_seeds3''.

\section{AUC vs L1 Selection}
\label{app:auc-l1}
% Three regimes from progress report §3.5: parsimony (gpt4/mistral/llama), convergence (gpt2), reversal (cohere). Demoted to appendix because the unconditional "L1 is more parsimonious than univariate AUC" claim doesn't hold across all 5 generators.

\section{Restricted Probe (Smaller-Model vs Ablate-Complement)}
\label{app:restricted}
% Three wide-gap generators: mpt (4.68±1.13pp), cohere (4.43±1.05pp), gpt2 (3.63±0.84pp); the other three are <1.1pp.
% Supporting evidence for the discussion's "both base generators + cohere have more dispersed signal" claim.

\section{Per-Domain Patching Breakdown}
\label{app:patching-perdomain}
% From progress report §3.3. Cohere uniformly high across domains (5–10\%); gpt2 strong on news/reddit/wiki (3–5\%) and weak on books/reviews (~2\%); chat-RLHF generators flat at 1–3.5\%.

\end{document}